\section{Non-IID Robust Aggregation}

\label{sec:aggregation}

\subsection{Ecological Non-IID Characterization}

We characterize the non-IID nature of wildlife monitoring data along three axes:

\begin{itemize}[leftmargin=*]
    \item \textbf{Species Distribution Skew:} The Jaccard similarity of species lists between sites ranges from 0.02 (Amazon vs.\ Serengeti) to 0.78 (adjacent sites in the same biome).
    \item \textbf{Observation Density Skew:} Camera trap densities range from 0.1 to 25 cameras per km$^2$, creating 250x variation in data volume per unit area.
    \item \textbf{Temporal Skew:} Monitoring periods range from 2 months to 8 years, with seasonal deployment patterns creating systematic temporal gaps.
\end{itemize}

\subsection{NIRA Algorithm}

Standard FedAvg computes a weighted average of client updates:

\begin{equation}
    \theta^{(t+1)} = \sum_{k=1}^{K} \frac{N_k}{N} \theta_k^{(t)}
\end{equation}

This performs poorly under extreme non-IID conditions because updates from sites with very different species assemblages may conflict, leading to catastrophic interference.

NIRA addresses this through hierarchical clustering-based aggregation:

\begin{enumerate}[leftmargin=*]
    \item \textbf{Ecological Clustering:} Sites are grouped into ecological clusters $\{C_1, \ldots, C_M\}$ based on species composition similarity, computed from the public metadata (species lists) without accessing raw data.
    \item \textbf{Intra-Cluster Aggregation:} Within each cluster, updates are averaged using standard FedAvg:
    \begin{equation}
        \theta_{C_m}^{(t+1)} = \sum_{k \in C_m} \frac{N_k}{\sum_{j \in C_m} N_j} \theta_k^{(t)}
    \end{equation}
    \item \textbf{Inter-Cluster Aggregation:} Cluster models are merged selectively, sharing only the components of the model that generalize across clusters (shared feature extractor) while maintaining cluster-specific classification heads:
    \begin{equation}
        \phi^{(t+1)} = \sum_{m=1}^{M} \frac{|C_m|}{K} \phi_{C_m}^{(t+1)} \qquad \psi_{C_m}^{(t+1)} = \text{local}
    \end{equation}
\end{enumerate}

\subsection{Dynamic Cluster Assignment}

Clusters are re-computed every $R$ rounds using a similarity metric based on gradient directions rather than raw data:

\begin{equation}
    \text{sim}(k, j) = \frac{\nabla_\phi \mathcal{L}_k \cdot \nabla_\phi \mathcal{L}_j}{\|\nabla_\phi \mathcal{L}_k\| \cdot \|\nabla_\phi \mathcal{L}_j\|}
\end{equation}

This allows clusters to adapt as models improve: sites that initially appear dissimilar may develop similar gradient directions as the shared representation matures.

\subsection{Rare Species Knowledge Transfer}

A key benefit of federation is improving detection of rare species through cross-site knowledge transfer.
We introduce a \emph{species embedding sharing} mechanism where sites with observations of a rare species contribute a species-specific embedding vector:

\begin{equation}
    \mathbf{e}_s^{\text{global}} = \frac{\sum_{k: s \in \mathcal{S}_k} N_{k,s} \cdot \mathbf{e}_s^{(k)}}{\sum_{k: s \in \mathcal{S}_k} N_{k,s}}
\end{equation}

\noindent where $\mathbf{e}_s^{(k)}$ is site $k$'s learned embedding for species $s$ and $N_{k,s}$ is the observation count.
These embeddings, being summary statistics rather than raw data, can be shared with minimal privacy risk (formally bounded by the DP guarantees on the gradient updates that produced them).

% ============================================================
% 7. Implementation
% ============================================================
